% META: {"id":"1SPE-TRIGO-RE-C2","chapitre":"1SPE-TRIGONOMETRIE","type_objet":"remediation","capacites_codes":["C2"],"status":"approved","origin":"generated","status_history":["generated","audited","accepted","approved"],"release_acceptance":"RELEASE_OWNER_BATCH_ACCEPTANCE","acceptance_closure_digest":"sha256:9b3ccf9a81c5520fbb7e03b7b2d3e7bf2057a02834b6d908bf3be5f49f3b553f","acceptance_receipt_digest":"sha256:4fad86f0282e0fa4e0419cca1ad6379ae7af5a280c04a4a90880f90a1c044242"}

%---------------------------------------------------------------
% FICHE DE REMEDIATION — C2 : Valeurs remarquables et symetries
%---------------------------------------------------------------

\begin{exercice}{1SPE-TRIGO-RE-C2-EX1}{1}{8}
\begin{enumerate}
  \item Compléter de memoire le tableau des valeurs remarquables pour $0$, $\dfrac{\pi}{6}$, $\dfrac{\pi}{4}$, $\dfrac{\pi}{3}$, $\dfrac{\pi}{2}$.
  \item Calculer $\cos\!\left(\dfrac{2\pi}{3}\right)$, $\sin\!\left(\dfrac{7\pi}{6}\right)$, $\cos\!\left(-\dfrac{\pi}{4}\right)$.
  \item Sachant que $\cos x = \dfrac{4}{5}$ et $x \in \left[0;\dfrac{\pi}{2}\right]$, calculer $\sin x$.
\end{enumerate}
\end{exercice}

% BEGIN-VERIFY
% from sympy import *
% assert cos(2*pi/3) == Rational(-1,2)
% assert sin(7*pi/6) == Rational(-1,2)
% assert cos(-pi/4) == sqrt(2)/2
% assert Rational(4,5)**2 + Rational(3,5)**2 == 1
% END-VERIFY

\begin{corrige}{1SPE-TRIGO-RE-C2-EX1}
\textbf{2.} $\cos\!\left(\dfrac{2\pi}{3}\right) = \cos(\pi-\dfrac{\pi}{3}) = -\cos\dfrac{\pi}{3} = -\dfrac{1}{2}$.

$\sin\!\left(\dfrac{7\pi}{6}\right) = \sin(\pi+\dfrac{\pi}{6}) = -\sin\dfrac{\pi}{6} = -\dfrac{1}{2}$.

$\cos(-\dfrac{\pi}{4}) = \cos\dfrac{\pi}{4} = \dfrac{\sqrt{2}}{2}$.

\textbf{3.} $\sin^2 x = 1 - \dfrac{16}{25} = \dfrac{9}{25}$, $\sin x = \dfrac{3}{5}$ (car $x \in [0;\dfrac{\pi}{2}]$).
\end{corrige}
